\section{Projective Observations and Recovery}\label{sec:projective-observation-boundary}

The hidden-operator comparison depends on the data actually observed. Full
Pl\"ucker data expose a subspace, even when no particular basis or matrix is
announced. This geometric recovery belongs to the HCP observation argument;
the commitment encoding and coding-theoretic consequences are developed in
Papers~2 and~3.

\subsection{Frames and Projective Coordinates}

Let \(V=\mathbb R^n\) have its fixed Euclidean inner product, and let
\(w_1,\ldots,w_r\) be an independent frame spanning \(W\le V\). Its exterior
coordinate is
\[
  c=w_1\wedge\cdots\wedge w_r\in\bigwedge^r V.
\]
An independent frame has nonzero exterior product. A different basis of the
same subspace changes \(c\) by the nonzero determinant of its change-of-basis
matrix. The projective class therefore depends only on \(W\):
\[
 c\sim c'\quad\Longleftrightarrow\quad
 \exists u\in\mathbb R^\times,\ c'=uc.
\]
The formal raw carrier is the scalar-orbit quotient of the whole exterior
power, including its zero orbit. A genuine frame determines a nonzero,
decomposable orbit. The conclusions below use such frames; they do not
assert that every point of the raw quotient is a Grassmannian point.

A related specialized reconstruction uses restricted Pl\"ucker coordinates
for maximal isotropic subspaces \cite{dingKim2026plucker}. The argument here
uses the full projective point of an arbitrary real subspace and the fixed
ambient metric; no orthogonal-matroid realization is assumed.

\begin{proposition}[Projective Pl\"ucker invariance]\label{prop:plucker-invariance}
  Two independent frames spanning the same finite-dimensional subspace have
  equal projective exterior coordinates.
\end{proposition}
\begin{proof}
  Both exterior products are nonzero elements of the one-dimensional top
  exterior power of that subspace. They differ by a nonzero scalar and hence
  have the same orbit.
\end{proof}

\subsection{Recovery from the Full Point}

\begin{proposition}[Recovery from the full projective point]
\label{prop:plucker-recovery}
  Let \(w_1,\ldots,w_r\) be an independent frame for
  \(W\subseteq\mathbb R^n\), and put
  \(c=w_1\wedge\cdots\wedge w_r\). Then
  \[
    W=\{x\in\mathbb R^n:x\wedge c=0\}.
  \]
  This reconstruction depends only on \([c]\). Consequently, two genuine
  rank-\(r\) frame witnesses have the same full projective point if and only
  if they present the same subspace. At the fixed ambient inner product,
  the full point determines \(P_W\) and \(I-P_W\).
\end{proposition}
\begin{proof}
  If \(x\in W\), the frame with \(x\) adjoined is dependent and its exterior
  product vanishes. If \(x\notin W\), the adjoined frame is independent and
  its exterior product is nonzero. Multiplying \(c\) by a nonzero scalar
  does not change this annihilator. The projector is the unique orthogonal
  projection onto the reconstructed subspace.
\end{proof}

In the raw quotient the zero orbit has the whole ambient space as its
annihilator. The frame-indexed statement avoids treating that orbit as an
arbitrary nonzero rank-\(r\) Pl\"ucker point.
If \(W=\ker\Delta_x\) is the operative harmonic subspace, full projective
data together with \(b\) determine
\[
  h=P_Wb.
\]
The operative-kernel identification and the fixed inner product are essential.
A coordinate for an unrelated subspace gives no such conclusion. The full
point does not determine the positive spectrum of the original \(\Delta_x\).
Nor does this exact-real reconstruction give an efficient finite byte
algorithm or inversion of a digest.

\subsection{Boundary with Protocol and Proximity Analysis}

A concrete protocol may map an original cochain payload to a finite received
word. Paper~2 owns that payload, its observation/encoding contract and the
stage evidence \cite{vanDerVlugt2026protocol}. Paper~3 owns the quantitative
source-to-code comparison, list bounds and MCA radii
\cite{vanDerVlugt2026proximity}. A comparison for its declared modular source
profile does not identify an arbitrary Euclidean or full-HCP residual with
code distance. In particular, the exact-component obstruction of
\cref{prop:exact-component-obstruction} must still be respected.

The present paper asserts no capacity-scale HCP-to-RS reduction merely from
those maps. Any such application needs a stated raw source carrier, its
valid locus, its metric and an independently justified far-preservation
comparison. The protocol and coding papers retain those hypotheses and
quantifier orders in their own statements; the detailed gate, serialization,
canonical-field and radius developments are not repeated here.
